% AA-17 | Algebra abstrakcyjna | audyt I
\chapter[Postać Smitha i klasyfikacja modułów]{Moduły nad dziedzinami ideałów głównych i postać Smitha}
\label{aa:ch:pid-modules-smith}

\section{Od prezentacji do klasyfikacji}
W Rozdziale~\ref{aa:ch:modules-introduction} pokazaliśmy, że każdy moduł
jest ilorazem modułu wolnego, a każdy moduł skończenie generowany jest
obrazem pewnego wolnego modułu skończonej rangi. W tym rozdziale pytamy, czy relacje między
nimi można uprościć tak, by otrzymać jednoznaczny opis modułu. Dla
ogólnych pierścieni odpowiedź jest przecząca. Dla \emph{dziedziny ideałów
głównych} sytuacja okazuje się jednak wyjątkowa: odpowiednią macierz
relacji można sprowadzić do postaci diagonalnej przez odwracalne
operacje na wierszach i kolumnach.

\begin{definicja}[Dziedzina ideałów głównych]
\label{aa:def:pid17}
Dziedziną ideałów głównych nazywamy niezerowy pierścień przemienny
$R$ z jedynką, bez dzielników zera, w którym każdy ideał ma postać
$(a)=aR$ dla pewnego $a\in R$. Element $u\in R$ jest \emph{jednostką},
jeśli istnieje $v\in R$ z $uv=1$. Elementy $a,b$ są \emph{stowarzyszone},
gdy $a=ub$ dla pewnej jednostki $u$. Dalej $R$ jest dziedziną ideałów
głównych, o ile nie zaznaczymy inaczej.
\end{definicja}

Przykłady obejmują $\mathbb Z$ oraz $K[t]$ dla dowolnego ciała $K$;
dla $K[t]$ stosujemy Twierdzenie~\ref{aa:tw:euclidean-pid}.
Pierścień $\mathbb Z[t]$ nie jest dziedziną ideałów głównych.
Istotnie, gdyby $(2,t)=(f)$, to wielomian $f$ dzieliłby
zarówno $2$, jak i $t$ w $\mathbb Z[t]$. Dzielnik $2$ jest
stały i równy $\pm1$ lub $\pm2$; ponieważ $2$ nie dzieli
$t$, otrzymalibyśmy $f=\pm1$ i $(2,t)=\mathbb Z[t]$.
Jest to sprzeczne z faktem, że homomorfizm
$\mathbb Z[t]\to\mathbb Z/2\mathbb Z$,
$g(t)\mapsto g(0)\bmod2$, zeruje $2$ i $t$, ale
nie zeruje $1$.

\begin{lemat}[Warunek stabilizacji rosnących ciągów ideałów]
\label{aa:lem:pid-acc17}
W dziedzinie ideałów głównych $R$ każdy rosnący ciąg ideałów
$I_1\subseteq I_2\subseteq\cdots$ od pewnego miejsca jest stały.
\end{lemat}
\begin{proof}
Suma mnogościowa $I=\bigcup_{n\ge1}I_n$ jest ideałem: dla $x\in I_m$
i $y\in I_n$ mamy $x-y\in I_{\max(m,n)}$, a mnożenie przez dowolny
element $R$ zachowuje tę samą składową. Ponieważ $R$ jest dziedziną
ideałów głównych, $I=(a)$. Generator $a$ należy do pewnego $I_N$.
Zatem $I=(a)\subseteq I_N\subseteq I$, czyli $I_N=I$ i wszystkie
$I_n$ dla $n\ge N$ są równe $I$. Nie użyliśmy funkcji euklidesowej;
argument działa w dowolnej dziedzinie ideałów głównych.
\end{proof}

\section{Podmoduły modułów wolnych i prezentacje skończone}
\begin{twierdzenie}[Podmoduł wolnego modułu skończonej rangi]
\label{aa:tw:pid-submodule-free}
Niech $R$ będzie dziedziną ideałów głównych, a $N\subseteq R^n$
podmodułem. Wtedy $N$ jest wolnym $R$-modułem o bazie zawierającej
co najwyżej $n$ elementów.
\end{twierdzenie}
\begin{proof}
Postępujemy indukcyjnie względem $n$. Dla $n=0$ moduł $R^0=0$,
a więc $N=0$ jest wolny z bazą pustą. Załóżmy twierdzenie dla $n-1$
i rozważmy projekcję $\pi:R^n\to R$ na pierwszą współrzędną.
Obraz $\pi(N)$ jest ideałem, ma zatem postać $aR$.

Jeżeli $a=0$, to $N\subseteq\{0\}\oplus R^{n-1}$ i teza wynika
z założenia indukcyjnego. Niech więc $a\ne0$. Wybierzmy $u\in N$
z pierwszą współrzędną $a$. Dla każdego $x\in N$ istnieje $r\in R$
takie, że $\pi(x)=ra$; wówczas $x-ru\in N\cap\ker\pi$.
Zatem $N=Ru+(N\cap\ker\pi)$. Suma jest prosta: jeśli
$ru\in\ker\pi$, to $ra=0$, a brak dzielników zera i $a\ne0$
dają $r=0$. Moduł $N\cap\ker\pi$ jest podmodułem $R^{n-1}$,
więc na mocy indukcji jest wolny z bazą co najwyżej $n-1$ elementów.
Dołączenie $u$ do tej bazy daje bazę $N$ długości co najwyżej $n$.
\end{proof}

\begin{wniosek}[Skończona prezentacja]
\label{aa:wn:pid-finite-presentation}
Niech $R$ będzie dziedziną ideałów głównych, a $M$
skończenie generowanym lewym $R$-modułem.
Wówczas $M$ ma prezentację
$R^m\xrightarrow{A}R^n\longrightarrow M\longrightarrow0$,
gdzie $m,n$ są skończone, a $A$ jest macierzą o współczynnikach w $R$.
\end{wniosek}
\begin{proof}
Wybierzmy generatory $m_1,\dots,m_n$ modułu $M$. Homomorfizm
$p:R^n\to M$ zadany przez $p(e_i)=m_i$ jest surjektywny.
Jądro $N=\ker p$ jest podmodułem $R^n$, a zatem jest wolne
z bazą $u_1,\dots,u_m$, gdzie $m\le n$. Odwzorowanie
$A:R^m\to R^n$ wysyłające $j$-ty wektor bazowy na $u_j$
ma obraz $N$. Otrzymujemy wymagany ciąg dokładny.
\end{proof}

\begin{uwaga}[Dlaczego założenie o dziedzinie jest istotne?]
Dla $R=K[\varepsilon]/(\varepsilon^2)$ ideał $(\varepsilon)$
nie jest wolnym $R$-modułem, choć jest podmodułem wolnego modułu $R$.
Rzeczywiście $\varepsilon(\varepsilon)=0$ i $\varepsilon\ne0$,
więc niezerowy skalar zeruje każdy jego element; w niezerowym
module wolnym nie zeruje każdego elementu bazowego. W pierścieniu
$R$ występują dzielniki zera, więc twierdzenie nie ma zastosowania.
\end{uwaga}

\section{Operacje elementarne i postać Smitha}
\begin{definicja}[Macierz relacji i operacje elementarne]
\label{aa:def:smith-operations}
Niech $A\in M_{m\times n}(R)$ wyznacza homomorfizm
$A:R^n\to R^m$. Dopuszczamy następujące operacje na wierszach
i kolumnach: zamianę miejscami, pomnożenie jednego wiersza lub
jednej kolumny przez jednostkę pierścienia oraz dodanie wielokrotności
jednego wiersza do innego (analogicznie dla kolumn).
Każda z tych operacji jest mnożeniem przez macierz odwracalną
z lewej lub z prawej strony. Ponadto w dowodzie postaci Smitha
będziemy wykonywać odwracalne zmiany baz na dwóch wybranych
wierszach lub kolumnach, opisane w
Lemacie~\ref{aa:lem:smith-bezout}. Są one również
mnożeniem przez macierz odwracalną, ale nad dowolną
dziedziną ideałów głównych nie zakładamy, że każdą taką
zmianę można rozłożyć na trzy powyższe typy prostych
operacji elementarnych. Rozróżniamy zatem
\emph{operacje elementarne} i ogólne \emph{odwracalne
zmiany baz}.
\end{definicja}

\begin{lemat}[Dwa elementy można zastąpić największym wspólnym dzielnikiem]
\label{aa:lem:smith-bezout}
Niech $R$ będzie dziedziną ideałów głównych, $a,b\in R$
i $(a,b)=(d)$, gdzie $d\ne0$.
Istnieje odwracalna macierz $U\in\operatorname{GL}_2(R)$ taka, że
$U(a,b)^T=(d,0)^T$.
\end{lemat}
\begin{proof}
Ponieważ $d$ generuje ideał $(a,b)$, istnieją $u,v\in R$
z $ua+vb=d$. Ponieważ $d$ dzieli $a$ i $b$, mamy $a=da'$,
$b=db'$ z $a',b'\in R$. Zapiszmy
\[
U=\begin{pmatrix}u&v\\-b'&a'\end{pmatrix}.
\]
Wyznacznik $U$ wynosi $ua'+vb'=(ua+vb)/d=1$, gdzie
skracanie przez niezerowy $d$ jest dozwolone, ponieważ $R$ nie
ma dzielników zera. Macierz $U$ jest zatem odwracalna: jej odwrotność to
$\begin{pmatrix}a'&-v\\b'&u\end{pmatrix}$.
Ponadto $U(a,b)^T=(d,-b'a+a'b)^T=(d,0)^T$.
Dla pary wpisanej w wiersz stosujemy macierz $U^T$ po prawej
stronie. Są to odwracalne zmiany baz, których poprawność
nie zależy od istnienia algorytmu dzielenia z resztą w $R$.
Nie utożsamiamy ich bez dowodu ze złożeniem operacji
elementarnych; nad dziedziną euklidesową taką redukcję
można natomiast realizować algorytmem Euklidesa.
\end{proof}

\begin{twierdzenie}[Postać normalna Smitha]
\label{aa:tw:smith-normal-form}
Niech $R$ będzie dziedziną ideałów głównych i
$A\in M_{m\times n}(R)$, gdzie $m,n\ge0$.
Istnieją odwracalne macierze $U\in\operatorname{GL}_m(R)$
i $V\in\operatorname{GL}_n(R)$ takie, że $UAV$ jest macierzą
prostokątną diagonalną, której niezerowe elementy diagonalne
$d_1,\dots,d_s$ spełniają
$d_1\mid d_2\mid\cdots\mid d_s$.
Pozostałe elementy są równe zeru. Jeśli $R$ jest dziedziną
euklidesową i dysponujemy efektywnym dzieleniem z resztą,
przekształcenia można wykonywać algorytmicznie.
Dla ogólnej dziedziny
ideałów głównych twierdzenie zapewnia istnienie
odwracalnych macierzy $U,V$, ale samo założenie o
ideałach głównych nie dostarcza efektywnej procedury
obliczania generatorów wszystkich występujących ideałów.
\end{twierdzenie}
\begin{proof}
Dla pustej macierzy lub macierzy zerowej twierdzenie jest
oczywiste. Załóżmy, że $A$ zawiera niezerowy współczynnik;
zamiana wierszy i kolumn ustawia go na pozycji $(1,1)$.
Nazwijmy ten element $a\ne0$.

Najpierw zbudujemy etap, w którym $a$ dzieli wszystkie
pozostałe współczynniki macierzy, a pierwsza kolumna i pierwszy
wiersz poza $a$ są zerowe. Jeśli pewien wpis $b$ pierwszej
kolumny nie jest podzielny przez $a$, zastosujmy do odpowiednich
dwóch wierszy Lemat~\ref{aa:lem:smith-bezout}. Powstały nowy
wpis w pozycji $(1,1)$ jest generatorem ideału $(a,b)$.
Ponieważ $a$ nie dzieli $b$, mamy ścisłe zawieranie
$(a)\subsetneq(a,b)$. Tak samo dla współczynnika pierwszego
wiersza, stosując operacje na dwóch kolumnach.

Po każdej takiej zmianie wpisu wiodącego ponownie sprawdzamy
pierwszy wiersz i kolumnę. Gdy $a$ dzieli już wszystkie ich
wpisy, przez odejmowanie
odpowiednich wielokrotności pierwszego wiersza od dalszych
wierszy zerujemy pierwszą kolumnę poniżej $a$. Następnie
odejmując wielokrotności pierwszej kolumny od pozostałych
kolumn zerujemy pierwszy wiersz poza $a$. Dostajemy
$\operatorname{diag}(a,B)$ dla pewnej macierzy $B$.

Jeśli $a$ dzieli wszystkie wpisy $B$, pierwszy etap jest zakończony.
Jeżeli natomiast w $B$ istnieje wpis $b$ niepodzielny przez $a$,
dodajmy jego wiersz do pierwszego wiersza całej macierzy.
Wpis na pozycji $(1,1)$ pozostaje równy $a$, ponieważ poniżej
niego są już zera; w pierwszym wierszu pojawia się za to
wpis $b$. Zastosowanie lematu do dwóch odpowiednich kolumn
zastępuje $a$ generatorem $(a,b)$ i ściśle powiększa ideał
wyznaczony przez wpis wiodący. Po tej operacji wracamy do
kontroli podzielności w pierwszym wierszu i kolumnie;
nie zakładamy, że wcześniejsze zera musiały się zachować.

W każdym kroku, w którym wpis wiodący się zmienia,
powstaje ściśle rosnący ideał główny $(a)$: jeśli
$a$ nie dzieli wybranego $b$, to $(a)\subsetneq(a,b)$.
Lemat~\ref{aa:lem:pid-acc17} wyklucza nieskończony
ciąg takich zmian. Przy ustalonym wpisie wiodącym
sprawdzamy skończenie wiele pozycji w pierwszym wierszu
i kolumnie, a następnie skończenie wiele wpisów bloku
$B$. Gdy któryś wpis nie jest podzielny przez $a$,
po najwyżej skończenie wielu działaniach otrzymujemy
nowy wpis wiodący generujący ściśle większy ideał.
Zatem po skończenie wielu operacjach otrzymujemy
macierz $\operatorname{diag}(d_1,B)$, gdzie $d_1\ne0$ dzieli
każdy współczynnik $B$.

Teraz postępujemy indukcyjnie z macierzą $B$ o wymiarach
$(m-1)\times(n-1)$. Operacje w jej wierszach i kolumnach
nie wpływają na pierwszy element diagonalny ani na zera poza
nim. Ponieważ wszystkie współczynniki $B$ są wielokrotnościami
$d_1$, każda późniejsza kombinacja liniowa tych współczynników
też jest wielokrotnością $d_1$. W szczególności uzyskany
niezerowy drugi element diagonalny $d_2$ jest podzielny przez
$d_1$. Powtarzając argument dostajemy żądaną postać.
Wszystkie operacje są odwracalne, stąd istnieją $U,V$.
Dla dziedziny euklidesowej wybór kombinacji i dzielenie
są dodatkowo efektywne.
\end{proof}

\begin{uwaga}[Ważne: dziedzina ideałów głównych nie musi być euklidesowa]
Dowód wykorzystuje stabilizację rosnących ciągów ideałów
oraz tożsamości Bézouta. W ogólnej dziedzinie ideałów
głównych stosujemy także odwracalne bloki macierzowe
$2\times2$; nie twierdzimy, że można je zawsze rozłożyć
na same operacje elementarne. W żadnym miejscu nie zakładamy
istnienia funkcji euklidesowej. Dlatego postać Smitha istnieje dla
każdej dziedziny ideałów głównych, nie tylko dla $\mathbb Z$
i $K[t]$.
\end{uwaga}

\section{Niezmienniki i jednoznaczność postaci Smitha}
\begin{definicja}[Ideały wyznaczone przez minory]
\label{aa:def:determinantal-ideals}
Niech $R$ będzie dziedziną ideałów głównych,
$A\in M_{m\times n}(R)$, gdzie $m,n\ge0$.
Dla $k\ge1$ oznaczamy przez $D_k(A)$ ideał
generowany przez wyznaczniki wszystkich podmacierzy kwadratowych
stopnia $k$. Kładziemy $D_0(A)=R$ oraz $D_k(A)=(0)$,
gdy $k>\min(m,n)$. Ideały $D_k(A)$ nazywamy
\emph{ideałami minorowymi} macierzy.
\end{definicja}
\begin{lemat}[Niezmienniczość ideałów minorowych]
\label{aa:lem:minor-invariant}
Niech $R$ będzie dziedziną ideałów głównych,
$A\in M_{m\times n}(R)$, $m,n\ge0$.
Dla dowolnych $U\in\operatorname{GL}_m(R)$ i
$V\in\operatorname{GL}_n(R)$ zachodzi $D_k(UAV)=D_k(A)$
dla każdego $k\ge0$.
\end{lemat}
\begin{proof}
Dla $k=0$ i $k>\min(m,n)$ teza wynika z definicji.
Niech zatem $1\le k\le\min(m,n)$. Wykorzystamy
rozwinięcie wyznacznika względem wyboru wierszy i kolumn.
Niech $I$ będzie zbiorem $k$ indeksów wierszy,
a $J$ zbiorem $k$ indeksów kolumn. Z wieloliniowości
wyznacznika i faktu, że wyznacznik z dwoma identycznymi
wierszami jest zerowy, otrzymujemy wzór
\[
 \det((UA)_{I,J})=
 \sum_{K:\,|K|=k}\det(U_{I,K})\det(A_{K,J}),
\]
gdzie $K$ przebiega $k$-elementowe zbiory indeksów
pośrednich. Aby zobaczyć ten wzór, rozpisujemy każdy wiersz
podmacierzy $(UA)_{I,J}$ jako kombinację wierszy $A$,
rozwijamy wyznacznik według wieloliniowości, odrzucamy
wyrazy z powtórzonym indeksem wiersza, a pozostałe
grupiujemy według rosnących ciągów indeksów $K$.
Współczynnik przy $\det(A_{K,J})$ jest wtedy dokładnie
$\det(U_{I,K})$, z uwzględnieniem znaków permutacji.

Każdy minor $UA$ jest więc kombinacją liniową minorów $A$,
stąd $D_k(UA)\subseteq D_k(A)$. Ponieważ $U$ jest
odwracalna, ten sam argument dla $U^{-1}$ daje
$D_k(A)\subseteq D_k(UA)$. Otrzymujemy równość.
Przeprowadzając analogiczne rozwinięcie dla kolumn
macierzy $AV$ i używając $V^{-1}$, dostajemy
$D_k(AV)=D_k(A)$. Złożenie obu równości kończy dowód.

\end{proof}

\begin{twierdzenie}[Jednoznaczność postaci Smitha]
\label{aa:tw:smith-uniqueness}
Niech $R$ będzie dziedziną ideałów głównych,
$A\in M_{m\times n}(R)$, $m,n\ge0$.
Niech $\operatorname{diag}(d_1,\dots,d_s,0,\dots)$ będzie
postacią Smitha macierzy $A$, gdzie $s\ge0$, $d_i\ne0$
i $d_i\mid d_{i+1}$ dla $1\le i<s$. Wówczas dla $1\le k\le s$
zachodzi $D_k(A)=(d_1\cdots d_k)$, natomiast
$D_k(A)=(0)$ dla $k>s$. Zatem liczba $s$ oraz
wszystkie $d_i$ są wyznaczone przez $A$ jednoznacznie
z dokładnością do stowarzyszenia.
\end{twierdzenie}
\begin{proof}
Z Lematu~\ref{aa:lem:minor-invariant} możemy liczyć
ideały minorowe dla postaci diagonalnej. Każdy niezerowy
minor stopnia $k$ jest, z dokładnością do znaku,
iloczynem $k$ różnych elementów diagonalnych.
Warunek $d_1\mid\dots\mid d_s$ sprawia, że
$d_1\cdots d_k$ dzieli każdy taki iloczyn, a sam
jest minorem z pierwszych $k$ wierszy i kolumn.
Dlatego ideał wszystkich minorów stopnia $k$ jest
$(d_1\cdots d_k)$. Dla $k>s$ wszystkie minory są zerowe.

Jeżeli dwie postacie Smitha mają współczynniki $d_i$
i $e_i$, z równości ideałów $D_1$ wynika, że $d_1$
i $e_1$ są stowarzyszone. Indukcyjnie porównujemy kolejne ideały
$(d_1\cdots d_k)=(e_1\cdots e_k)$. Generatory
równego niezerowego ideału są stowarzyszone, więc
$d_1\cdots d_k=u(e_1\cdots e_k)$ dla pewnej jednostki $u$.
Z założenia indukcyjnego
$d_1\cdots d_{k-1}=v(e_1\cdots e_{k-1})$
dla pewnej jednostki $v$. Skracając niezerowy iloczyn
$e_1\cdots e_{k-1}$ w dziedzinie całkowitości,
dostajemy $d_k=uv^{-1}e_k$. Zatem również $d_k$
i $e_k$ są stowarzyszone. Liczba
niezerowych współczynników to największe $k$
z $D_k(A)\ne(0)$, więc również jest niezmiennikiem.
\end{proof}

\begin{przyklad}[Rachunek postaci Smitha nad liczbami całkowitymi]
\label{aa:prz:smith-integer-steps}
Niech $A=\begin{pmatrix}4&6\\10&14\end{pmatrix}$.
Odejmując dwukrotność pierwszego wiersza od drugiego,
a następnie zamieniając wiersze, otrzymujemy
$\begin{pmatrix}2&2\\4&6\end{pmatrix}$.
Odejmując dwukrotność pierwszego wiersza od drugiego,
dostajemy $\begin{pmatrix}2&2\\0&2\end{pmatrix}$.
Na koniec odejmujemy pierwszą kolumnę od drugiej:
$\operatorname{diag}(2,2)$. W szczególności
$\mathbb Z^2/A\mathbb Z^2\cong C_2\oplus C_2$.
Bez wykonania operacji wynik kontrolujemy przez
$\gcd(4,6,10,14)=2$ i $|\det A|=4$:
pierwszy współczynnik Smitha ma wartość bezwzględną $2$,
a iloczyn dwóch współczynników ma wartość bezwzględną $4$.
\end{przyklad}

\begin{przyklad}[Dwa równania relacji i iloraz]
\label{aa:prz:smith-cokernel-90}
Dla $A=\begin{pmatrix}6&9\\0&15\end{pmatrix}$
mamy $D_1(A)=(3)$, $D_2(A)=(90)$. Dlatego
postać Smitha ma współczynniki stowarzyszone z $3$ i $30$.
Moduł o dwóch generatorach $u,v$ i relacjach
$6u=0$, $9u+15v=0$ jest izomorficzny z
$\mathbb Z/3\mathbb Z\oplus\mathbb Z/30\mathbb Z$.
Iloraz ma $90$ elementów, ponieważ ma strukturę
$C_3\oplus C_{30}$. Nie należy utożsamiać liczby
elementów z liczbą generatorów: wystarczają dwa.
W rzeczywistości jeden nie wystarczy, gdyż
$M/3M\cong(\mathbb Z/3\mathbb Z)^2$ ma wymiar dwa
nad $\mathbb Z/3\mathbb Z$.
\end{przyklad}

\section{Klasyfikacja skończenie generowanych modułów}
\begin{twierdzenie}[Twierdzenie o strukturze modułów nad dziedziną ideałów głównych]
\label{aa:tw:pid-structure}
Niech $R$ będzie dziedziną ideałów głównych, a $M$
skończenie generowanym $R$-modułem. Istnieją liczby całkowite
$r,t\ge0$ oraz niezerowe, nieodwracalne elementy
$d_1,\dots,d_t\in R$ spełniające $d_1\mid\cdots\mid d_t$
(dla $t=0$ rodzina jest pusta) takie, że
\[
 M\cong R^r\oplus\bigoplus_{i=1}^t R/(d_i).
\]
Liczba $r$ jest wyznaczona jednoznacznie, a ciąg
$d_1,\dots,d_t$ jest jednoznaczny z dokładnością do
pomnożenia każdego elementu przez jednostkę $R$.
\end{twierdzenie}
\begin{proof}
Najpierw dowiedziemy istnienia. Z
Wniosku~\ref{aa:wn:pid-finite-presentation} mamy
$M\cong\operatorname{coker}(A)$ dla pewnego
$A:R^m\to R^n$. Z Twierdzenia~\ref{aa:tw:smith-normal-form}
otrzymujemy $UAV=D$, gdzie $U$ i $V$ są odwracalne.
Ponieważ $V$ jest automorfizmem modułu źródłowego,
$\operatorname{im}(AV)=\operatorname{im}A$.
Ponieważ $U$ jest automorfizmem modułu docelowego,
indukuje izomorfizm $R^n/\operatorname{im}A\cong
R^n/\operatorname{im}D$. Dla macierzy diagonalnej
iloraz rozkłada się współrzędnie na
$R/(d_1)\oplus\cdots\oplus R/(d_s)\oplus R^{n-s}$.
Jeśli $d_i$ jest jednostką, odpowiedni czynnik jest zerowy;
pozostałe niezerowe współczynniki zachowują warunek
podzielności. To dowodzi istnienia wskazanego rozkładu.

Dla jednoznaczności najpierw oznaczmy przez $T(M)$ zbiór
elementów $m\in M$, dla których $am=0$ dla pewnego
$a\in R\setminus\{0\}$. Jeżeli $am=0$ i $bn=0$ dla niezerowych $a,b\in R$,
to $ab(m-n)=0$, przy czym $ab\ne0$. Dla dowolnego
$c\in R$ mamy ponadto $a(cm)=c(am)=0$, dlatego
$T(M)$ jest podmodułem. Każdy czynnik $R/(d_i)$
jest zerowany przez $d_i$, a cała skończona suma tych
czynników przez iloczyn $d_1\cdots d_t$ (dla $t=0$
przez $1$). Natomiast niezerowy wektor w $R^r$ nie
może zostać wyzerowany przez niezerowy skalar, ponieważ
$R$ nie ma dzielników zera. Stąd w podanym rozkładzie
$T(M)\cong\bigoplus_i R/(d_i)$ i $M/T(M)\cong R^r$.
Dwie rangi wolnych modułów nad $R$ są równe także wtedy,
gdy $R$ nie jest ciałem. Niech $K=\operatorname{Frac}(R)$
będzie ciałem ułamków dziedziny $R$. Izomorfizm
$R^r\cong R^s$ jest reprezentowany przez macierz
$A\in M_{s\times r}(R)$, a jego odwrotność przez
$B\in M_{r\times s}(R)$, z $BA=I_r$ i $AB=I_s$.
Traktując współczynniki jako elementy $K$, dostajemy
wzajemnie odwrotne $K$-liniowe odwzorowania
$K^r\to K^s$ i $K^s\to K^r$. Z równości
wymiarów izomorficznych przestrzeni wektorowych
wynika $r=s$. Argument ten działa dla każdej
dziedziny całkowitości i nie wymaga wcześniejszej
teorii iloczynu tensorowego.

Pozostaje jednoznaczność części torsyjnej. Z
Twierdzenia~\ref{aa:tw:pid-ufd} każdy niezerowy
nieodwracalny element $R$ rozkłada się jednoznacznie
na czynniki nierozkładalne z dokładnością do jednostek
i kolejności. Niech $p\in R$ będzie elementem
nierozkładalnym. Ideał $(p)$ jest maksymalny:
jeśli $(p)\subseteq(a)\subsetneq R$, to $a$ dzieli
$p$, a ponieważ $a$ nie jest jednostką, z nierozkładalności
$p$ wynika, że $a$ jest stowarzyszony z $p$.
Stąd $R/(p)$ jest ciałem.
Dla dowolnego modułu torsyjnego $T$ oznaczmy
$T[p^k]=\{x\in T:p^kx=0\}$ dla $k\ge1$ oraz
$T[p^0]=0$. Iloraz $T[p^k]/T[p^{k-1}]$ jest
przestrzenią wektorową nad ciałem $R/(p)$:
mnożenie przez $p$ wysyła $T[p^k]$ do $T[p^{k-1}]$,
więc działa na ilorazie jako zero.
Dla pojedynczego czynnika $L=R/(p^e)$, gdzie $e\ge1$,
zbiór $L[p^k]$ jest, dla $1\le k\le e$, podmodułem
generowanym przez klasę $p^{e-k}$. Istotnie,
$p^k(x+(p^e))=0$ dokładnie wtedy, gdy $p^e\mid p^kx$;
przez skracanie w dziedzinie otrzymujemy
$p^{e-k}\mid x$. Dla $k>e$ mamy $L[p^k]=L$.
Gdy $1\le k\le e$, odwzorowanie
\[
 R/(p)\longrightarrow L[p^k]/L[p^{k-1}],\qquad
 c+(p)\longmapsto cp^{e-k}+L[p^{k-1}]
\]
jest dobrze określone: zastąpienie $c$ przez $c+pa$
zmienia reprezentanta o $ap^{e-k+1}$, który należy do
$L[p^{k-1}]$. Odwzorowanie jest na z opisu $L[p^k]$
i ma jądro zerowe, ponieważ $cp^{e-k}$ należy do
$L[p^{k-1}]$ dokładnie wtedy, gdy $p\mid c$.
Jest więc izomorfizmem $R/(p)$-przestrzeni wektorowych.
Wobec tego rozważany iloraz ma wymiar jeden dla
$1\le k\le e$ oraz zero dla $k>e$.
Każdy $R/(d_i)$ rozkłada się na sumę składników
$R/(p^{e_{i,p}})$ po rozłożeniu $d_i$ na potęgi
parami niestowarzyszonych elementów nierozkładalnych
i zastosowaniu chińskiego twierdzenia o resztach
(Twierdzenie~\ref{aa:tw:crt-ring}): ideały
$(p^e)$ i $(q^f)$ dla niestowarzyszonych
nierozkładalnych $p,q$ są względnie pierwsze.
Istotnie, ich suma ma postać $(g)$, a $g$ dzieli
$p^e$ i $q^f$, więc na mocy jednoznaczności
rozkładu $g$ jest jednostką. Zatem
$(p^e)+(q^f)=R$. Jeżeli $p\nmid d_i$, to z jednoznaczności rozkładu na
czynniki mamy $(p^k,d_i)=R$. Istnieją zatem $u,v\in R$
takie, że $up^k+vd_i=1$. Jeśli $p^k$ zeruje klasę
$x+(d_i)$, to po pomnożeniu tej tożsamości przez $x$
widzimy, że $x\in(d_i)$; stąd $R/(d_i)$ nie ma
niezerowych elementów zerowanych przez $p^k$.
W konsekwencji liczba
\[
 h_{p,k}=\dim_{R/(p)}\bigl(T[p^k]/T[p^{k-1}]\bigr)
\]
jest równa liczbie tych indeksów $i$, dla których
$e_{i,p}\ge k$. Wartości wszystkich $h_{p,k}$ są
niezmiennikami modułu, ponieważ $T[p^k]$ definiujemy
bez użycia wybranego rozkładu, a każdy izomorfizm
modułów zachowuje te podmoduły i wywołuje izomorfizmy
ich kolejnych ilorazów. Pozwalają one odtworzyć liczbę
wystąpień wykładnika dokładnie $k$ jako
$h_{p,k}-h_{p,k+1}$.

Dla każdego $p$ odzyskujemy więc multizbiór wykładników
$e_{i,p}>0$. Warunek $d_1\mid\cdots\mid d_t$ oznacza,
że dla każdego $p$ wykładniki w kolejnych $d_i$
są niemalejące; zerowe wykładniki dopełniamy od lewej.
Jeśli $T(M)=0$, to $t=0$ i nie ma czynników torsyjnych.
W przeciwnym przypadku $t$ jest największą liczbą
dodatnich wykładników wśród różnych $p$: ponieważ
$d_1$ nie jest jednostką, ma przynajmniej jeden
czynnik nierozkładalny $p$, który dzieli wszystkie
$d_i$; dla tego $p$ wszystkie $t$ wykładników są dodatnie.
Tak ustalone ciągi wykładników określają wszystkie $d_i$
z dokładnością do jednostek pierścienia. To kończy dowód.
\end{proof}

\begin{wniosek}[Klasyfikacja skończenie generowanych grup abelowych]
\label{aa:wn:abelian-fg-structure}
Każda skończenie generowana grupa abelowa $A$ jest izomorficzna z
$\mathbb Z^r\oplus C_{n_1}\oplus\cdots\oplus C_{n_t}$,
gdzie $n_i\ge2$ i $n_1\mid\cdots\mid n_t$.
Liczba $r$ i wartości $n_i$ są jednoznaczne.
\end{wniosek}
\begin{proof}
Grupy abelowe są dokładnie modułami nad $\mathbb Z$.
Ponieważ $\mathbb Z$ jest dziedziną ideałów głównych,
stosujemy Twierdzenie~\ref{aa:tw:pid-structure}.
W $\mathbb Z$ jednostkami są $1$ i $-1$,
więc współczynniki $d_i$ można wybrać dodatnie.
Otrzymane czynniki $\mathbb Z/(d_i)$ to grupy cykliczne
$C_{d_i}$, co daje tezę.
\end{proof}

\begin{przyklad}[Grupy abelowe rzędu $72$]
\label{aa:prz:abelian-order72}
Ponieważ $72=2^3\cdot3^2$, najpierw rozkładamy grupę
na jej części $2$-pierwotną i $3$-pierwotną.
Dla liczby $8$ możliwe rozkłady potęgowo-pierwotne to
$C_8$, $C_4\oplus C_2$ i $C_2^3$; odpowiadają one
podziałom wykładnika $3$: $3$, $2+1$ i $1+1+1$.
Dla liczby $9$ są dwie możliwości: $C_9$ oraz $C_3^2$.
Wobec jednoznaczności klasyfikacji istnieje dokładnie
$3\cdot2=6$ klas izomorfizmu grup abelowych rzędu $72$.
W postaci czynników niezmienniczych reprezentują je
$C_{72}$, $C_3\oplus C_{24}$, $C_2\oplus C_{36}$,
$C_6\oplus C_{12}$, $C_2\oplus C_2\oplus C_{18}$ oraz
$C_2\oplus C_6\oplus C_6$. Każdy z tych rozkładów
spełnia wymagany warunek podzielności.
Przykładowo z chińskiego twierdzenia o resztach mamy
$C_4\oplus C_2\oplus C_9\cong C_2\oplus C_{36}$.
Ta grupa ma cztery elementy o rzędzie dzielącym $2$,
podczas gdy $C_{72}$ ma tylko dwa, więc grupy
nie są izomorficzne.
\end{przyklad}

\begin{przyklad}[Dlaczego hipotezy twierdzenia są potrzebne?]
Moduł $\mathbb Z/2\mathbb Z$ nie jest wolny nad $\mathbb Z$,
ale jest sumą czynników cyklicznych zgodnie z twierdzeniem.
Nad pierścieniem $S=K[t_1,t_2]$, który nie jest dziedziną
ideałów głównych, ideał $I=(t_1,t_2)$ jest skończenie generowanym
modułem, ale nie jest wolny. Rzeczywiście, gdyby $I$ miał bazę zawierającą
przynajmniej dwa elementy $u,v$, to niezerowa relacja
$vu-uv=0$ (ze współczynnikami $v,-u\in S$)
przeczyłaby liniowej niezależności tej bazy.
Niezerowy wolny podmoduł $I\subset S$ musiałby zatem
mieć rangę jeden, czyli być generowany przez
jeden element jako ideał. Tymczasem wspólnymi dzielnikami
$t_1$ oraz $t_2$ w $S$ są tylko jednostki, a
$1\notin(t_1,t_2)$, co widać po podstawieniu
$t_1=t_2=0$. Ideał $I$ nie jest zatem główny.
Nie wolno bez dodatkowych założeń przenosić twierdzenia
o strukturze na wszystkie pierścienie przemienne.
\end{przyklad}

\section{Dalszy związek z algebrą liniową}
Dla $K$ będącego ciałem pierścień $K[t]$ jest dziedziną
euklidesową, a więc dziedziną ideałów głównych. Jeśli
$V$ jest skończenie wymiarową przestrzenią wektorową nad $K$
i $T:V\to V$ jest operatorem liniowym, możemy traktować $V$
jako $K[t]$-moduł przez $t\cdot v=T(v)$.
Wówczas postać Smitha i twierdzenie o strukturze modułów
prowadzą do postaci kanonicznej Frobeniusa oraz do teorii
czynników niezmienniczych operatora. Pełne dowody tych
rezultatów umieścimy w późniejszym skrypcie algebry liniowej.
